\subsection{Limitations of the Thermal Guard}
\label{sec:guard-limit}
The goal is to enable lightweight multi-frame composition under parallel capture whenever the draft sequence can finish within its deadline.

\smallskip\noindent\textbf{Production thermal guard.}
Samsung's overheat level is an ordinal thermal indicator on a 0--6 scale, where larger values denote stronger throttling.
During pre-release validation, throttling at level 4 delayed a single-frame draft stage and caused a Capture Timeout.
The failure recurred despite optimization, so we introduced a thermal guard that skips optional stages at level 4 or higher.
Because the draft sequence is implemented in framework code shared across device models, the same cutoff applies to all supported models.

The overheat level does not indicate a capture's waiting time or remaining draft work.
The guard may therefore skip optional stages even when enough time remains to complete the draft sequence.
Ambient heat, charging, or preceding device workloads can also raise the overheat level to 4 during ordinary use.
If a session starts at this level, the guard skips optional stages even for its first capture.

\smallskip\noindent\textbf{Motivating experiment.}
To examine when the extended draft sequence can complete within its deadline, we evaluated two Portrait configurations.
The optional stages comprise the lightweight multi-frame draft stage, denoted by \(M\), and the deployed single-frame stages, collectively denoted by \(S\).
We compared \(M\) alone with \(M+S\), both followed by the same mandatory encoding and saving stages.
We selected the Galaxy S26 Ultra because it lacks the additional parallel-capture restrictions used by other flagship variants.

With the thermal guard bypassed, we ran ten 30-capture trials for each combination of starting overheat level (0--6), resolution (12MP or 24MP), memory condition, and draft sequence configuration.
An Android Debug Bridge (ADB) loop requested each subsequent capture as soon as the application permitted it.
The motivating level-4 failure occurred near the 20th consecutive capture with 12 draft sequences pending.
We therefore extended each trial to 30 captures to exercise the draft pipeline beyond that observed failure point.
For the memory-pressure condition, we ran the camera alongside AnTuTu Benchmark~11.1.4's memory-stress workload.
Table~\ref{tab:timeout_index} reports the earliest timeout capture index across the ten trials for each condition.

% Notes: docs/exhibits.md#tab_timeout_index
\providecolor{tmoS1}{RGB}{252,227,220} %
\providecolor{tmoS2}{RGB}{239,160,140} %
\providecolor{tmoS3}{RGB}{201,80,63}  %
\providecolor{tmoS4}{RGB}{143,36,32}  %

\newcolumntype{H}{>{\centering\arraybackslash}m{1.20cm}}
\newcolumntype{P}{>{\centering\arraybackslash}m{0.67cm}}

\providecommand{\tmoFirst}[2]{\cellcolor{#1}\makebox[\linewidth][r]{#2\hspace{1pt}}}
\providecommand{\tmoFirstW}[2]{\cellcolor{#1}\makebox[\linewidth][r]{\textcolor{white}{\textbf{#2}}\hspace{1pt}}}
\providecommand{\tmoNone}{\makebox[\linewidth][r]{--\hspace{1pt}}}

\begin{table}[t]
 \centering
 \caption{Earliest Timeout onset across capture conditions.}
 \label{tab:timeout_index}
 {\scriptsize\setlength{\tabcolsep}{1.3pt}%
 \begin{tabular}{@{}HPP|PP||PP|PP@{}}
 \toprule
 \multirow[c]{3}{=}[-10pt]{\makecell[c]{\textbf{Starting}\\\textbf{overheat}\\\textbf{level}}} &
 \multicolumn{4}{c||}{\raisebox{-3pt}{\textbf{Normal capture}}} &
 \multicolumn{4}{c@{}}{\raisebox{-3pt}{\textbf{Memory-pressure capture}}} \\[6pt] \cmidrule(lr){2-5}\cmidrule(l){6-9}
 &
 \multicolumn{2}{c|}{\raisebox{-3pt}{\textbf{12MP}}} &
 \multicolumn{2}{c||}{\raisebox{-3pt}{\textbf{24MP}}} &
 \multicolumn{2}{c|}{\raisebox{-3pt}{\textbf{12MP}}} &
 \multicolumn{2}{c@{}}{\raisebox{-3pt}{\textbf{24MP}}} \\[6pt] \cmidrule(lr){2-3}\cmidrule(lr){4-5}\cmidrule(lr){6-7}\cmidrule(l){8-9}
 & \raisebox{-2pt}{\textbf{M{+}S}} & \raisebox{-2pt}{\textbf{M}} &
 \raisebox{-2pt}{\textbf{M{+}S}} & \raisebox{-2pt}{\textbf{M}} &
 \raisebox{-2pt}{\textbf{M{+}S}} & \raisebox{-2pt}{\textbf{M}} &
 \raisebox{-2pt}{\textbf{M{+}S}} & \raisebox{-2pt}{\textbf{M}} \\[4pt]
 \midrule
 Lv0 &
 \tmoNone & \tmoNone & \tmoFirst{tmoS1}{29} & \tmoNone &
 \tmoNone & \tmoNone & \tmoFirst{tmoS1}{26} & \tmoNone \\
 Lv1 &
 \tmoNone & \tmoNone & \tmoFirst{tmoS1}{29} & \tmoNone &
 \tmoFirst{tmoS1}{24} & \tmoNone & \tmoFirst{tmoS2}{19} & \tmoNone \\
 Lv2 &
 \tmoFirst{tmoS1}{24} & \tmoNone & \tmoFirst{tmoS2}{18} & \tmoNone &
 \tmoFirst{tmoS2}{15} & \tmoFirst{tmoS1}{29} &
 \tmoFirst{tmoS2}{14} & \tmoFirst{tmoS1}{25} \\
 Lv3 &
 \tmoFirst{tmoS2}{13} & \tmoFirst{tmoS1}{21} &
 \tmoFirst{tmoS2}{12} & \tmoFirst{tmoS1}{20} &
 \tmoFirstW{tmoS3}{8} & \tmoFirst{tmoS2}{19} &
 \tmoFirstW{tmoS4}{5} & \tmoFirst{tmoS2}{13} \\
 \midrule
 Lv4 &
 \tmoFirstW{tmoS3}{8} & \tmoFirst{tmoS2}{16} &
 \tmoFirstW{tmoS3}{7} & \tmoFirst{tmoS2}{13} &
 \tmoFirstW{tmoS3}{6} & \tmoFirst{tmoS2}{13} &
 \tmoFirstW{tmoS4}{3} & \tmoFirstW{tmoS3}{7} \\
 Lv5 &
 \tmoFirstW{tmoS3}{8} & \tmoFirst{tmoS2}{12} &
 \tmoFirstW{tmoS3}{8} & \tmoFirst{tmoS2}{12} &
 \tmoFirstW{tmoS4}{4} & \tmoFirstW{tmoS3}{9} &
 \tmoFirstW{tmoS4}{4} & \tmoFirstW{tmoS3}{9} \\
 Lv6 &
 \tmoFirstW{tmoS3}{7} & \tmoFirst{tmoS2}{11} &
 \tmoFirstW{tmoS3}{7} & \tmoFirst{tmoS2}{11} &
 \tmoFirstW{tmoS4}{3} & \tmoFirstW{tmoS3}{6} &
 \tmoFirstW{tmoS4}{3} & \tmoFirstW{tmoS3}{6} \\
 \bottomrule
 \end{tabular}%
 }
 \par\vspace{3pt}
 {\scriptsize\begin{minipage}{0.933\columnwidth}
 \raggedright\setlength{\fboxsep}{1.8pt}%
 Earliest timeout onset:~%
 \colorbox{tmoS1}{\strut\,${\geq}\,20$\,}\colorbox{tmoS2}{\strut\,10--19\,}\colorbox{tmoS3}{\strut\,\textcolor{white}{6--9}\,}\colorbox{tmoS4}{\strut\,\textcolor{white}{${\leq}\,5$}\,}\par
 \end{minipage}}
\end{table}

With \(M+S\), timeouts occurred across the full range of starting overheat levels.
Retaining \(M\) while omitting \(S\) delayed the earliest observed timeout in every affected condition or yielded no timeout within 30 captures.
Under normal 12MP capture starting at level 4, however, \(M\) alone still timed out: the earliest timeout moved from capture 8 with \(M+S\) to capture 16.
Starting lower did not necessarily postpone failure: the earliest \(M+S\) timeout occurred at capture 5 under memory-pressure 24MP capture starting at level 3.
All ten normal 12MP \(M+S\) trials starting at level 4 completed the first seven captures, whereas the guard would skip \(M\) at that starting level.

\smallskip\noindent\textbf{Need for runtime control.}
The thermal cutoff does not capture how earlier draft sequences consume a later capture's time budget.
Skipping optional stages shortens that capture's processing but cannot recover time already spent waiting.
Executing those stages can also lengthen the wait for subsequent captures.
Delaying new requests can limit queue growth, but lengthens the shot-to-shot interval.
These constraints motivate joint runtime control of optional stages and capture arrivals.
